% The four parts of the Day 2 lab. The steps follow Labs/day02/README.md.
%
% A value that only a board can give is written as "..." and the students
% measure it. The deck does not show a result that the students must predict.

% ------------------------------------------------------------------- Part A
\begin{frame}[fragile]{Part A: write the MicroPython firmware (40 min)}
  \small
  \begin{enumerate}
    \item Download the firmware file and check it.
    \item Start bootloader mode: disconnect the USB cable, hold the \code{B}
          (boot) button of the XIAO, connect the cable, release the button.
    \item Erase the flash and write the firmware.
    \item Press the \code{RST} button of the expansion board.
  \end{enumerate}
\begin{lstlisting}[style=shell]
bash ../hardware/HW-01/get_firmware.sh
python3 -m esptool --chip esp32s3 --port PORT erase_flash
python3 -m esptool --chip esp32s3 --port PORT write_flash 0 \
  ../hardware/HW-01/downloads/SEEED_XIAO_ESP32S3-20260824-v1.29.0.bin
\end{lstlisting}
  \begin{itemize}
    \item \code{PORT} is the serial port of the board: \code{/dev/ttyACM0}
          (Linux), \code{/dev/cu.usbmodem101} (macOS), or \code{COM4}
          (Windows).
    \item The port name can change after the reset. The command
          \code{mpremote devs} lists the ports.
    \item The boot button is on the XIAO, not on the expansion board. It is
          very small.
  \end{itemize}
\end{frame}

\begin{frame}[fragile]{Part A: the REPL and a GPIO pin}
  \begin{columns}[T]
    \begin{column}{0.50\textwidth}
\begin{lstlisting}[style=shell]
mpremote connect PORT repl
\end{lstlisting}
\begin{lstlisting}[style=python]
>>> import sys, machine, gc
>>> sys.implementation
>>> machine.freq()
>>> gc.mem_free()

>>> from machine import Pin
>>> led = Pin(21, Pin.OUT)
>>> led.value(0)       # LED on
>>> led.value(1)       # LED off
\end{lstlisting}
    \end{column}
    \begin{column}{0.47\textwidth}
      \small
      \begin{itemize}
        \item The REPL is the prompt of MicroPython. You enter one line and
              you see the result.
        \item \textbf{You record:} the MicroPython version, the clock
              frequency in Hz, and the free memory in bytes.
        \item \textbf{You see:} the LED of the XIAO is on when the pin is
              low.
        \item Press \code{Ctrl-]} to close the REPL.
      \end{itemize}
    \end{column}
  \end{columns}
  \vskip 0.4em
  \keypoint{Part 1 of the lecture gave the clock and the memory of the
    board. The REPL shows the two values that your program has.}
\end{frame}

\begin{frame}[fragile]{Part A: timer, interrupt, and I2C scan}
\begin{lstlisting}[style=shell]
mpremote connect PORT run board/blink_timer.py
\end{lstlisting}
\begin{lstlisting}
The LED changes two times each second. Press the boot button.
ticks = ..., presses = ...
\end{lstlisting}
  \small
  \begin{itemize}
    \item A hardware timer calls a handler two times each second. The
          handler changes the LED and counts.
    \item A pin interrupt counts each press of the boot button.
    \item \textbf{You record:} the number of ticks after 10 s.
  \end{itemize}
\begin{lstlisting}[style=shell]
mpremote connect PORT run board/i2c_scan.py
\end{lstlisting}
\begin{lstlisting}
Found 2 device(s)
  0x3C  OLED display (SSD1306)
  0x6A  IMU (LSM6DS3TR-C)
\end{lstlisting}
  \begin{itemize}
    \item \textbf{You record:} the two addresses.
  \end{itemize}
  \source{The two outputs show the form that the scripts print.}
\end{frame}

% ------------------------------------------------------------------- Part B
\begin{frame}[fragile]{Part B: read the IMU (50 min)}
\begin{lstlisting}[style=shell]
mpremote connect PORT fs cp board/lsm6ds3.py :lsm6ds3.py
mpremote connect PORT repl
\end{lstlisting}
\begin{lstlisting}[style=python]
>>> from machine import I2C, Pin
>>> from lsm6ds3 import LSM6DS3
>>> imu = LSM6DS3(I2C(0, sda=Pin(5), scl=Pin(6), freq=400000))
>>> hex(imu.chip_id)       # the identity register: 0x6a
>>> imu.read()             # (ax, ay, az, gx, gy, gz)
\end{lstlisting}
  \small
  \begin{itemize}
    \item The first command copies the driver to the file system of the
          board. The file stays there after a reset.
    \item \code{imu.read()} gives the acceleration in $g$ and the angular
          rate in degrees each second.
    \item \textbf{You see:} with the kit flat on the table, one acceleration
          axis is near 1.0 $g$. The other two are near 0.
    \item \textbf{You record:} the six values.
  \end{itemize}
  \keypoint{The driver has about 130 lines. Read it: it uses the registers
    of Part 1 of the lecture.}
\end{frame}

\begin{frame}[fragile]{Part B: predict the rate, then measure it}
  \small
  The first version of \code{board/imu\_stream.py} reads the sensor, prints
  one line, and then waits 20 ms.
  \begin{exerciseblock}
    \textbf{Prediction.} Is the real rate equal to 50 Hz, lower, or higher?
    By how much? Write your prediction in the report before you measure.
  \end{exerciseblock}
\begin{lstlisting}[style=shell]
mpremote connect PORT fs cp board/imu_stream.py :main.py
mpremote connect PORT reset
python3 host/check_rate.py --port PORT --seconds 20 --save rate_wait.csv
\end{lstlisting}
\begin{lstlisting}
Samples:            ...
Lost lines:         ...
Mean rate:          ... Hz (nominal 50.00 Hz)
Period min / max:   ... us / ... us
RESULT: ...
\end{lstlisting}
  \begin{itemize}
    \item A script with the name \code{main.py} starts at each power-on.
    \item \textbf{You record:} the mean rate and the result.
  \end{itemize}
  \source{The output shows the form of the lines.}
\end{frame}

\begin{frame}[fragile]{Part B: task B1, a deadline for each sample}
  \begin{columns}[T]
    \begin{column}{0.47\textwidth}
      \small
      Change the end of the loop in \code{board/imu\_stream.py}:
      \begin{enumerate}
        \item Start \code{deadline} with \code{time.ticks\_us()} before
              the loop.
        \item Add \code{PERIOD\_US} with \code{time.ticks\_add()}.
        \item Get the time to the deadline with \code{time.ticks\_diff()}.
        \item Sleep only for this time.
        \item If the time is negative, start a new time base.
      \end{enumerate}
    \end{column}
    \begin{column}{0.50\textwidth}
\begin{lstlisting}[style=python]
# The first version: replace this line.
time.sleep_us(PERIOD_US)
\end{lstlisting}
\begin{lstlisting}[style=shell]
mpremote connect PORT fs cp \
  board/imu_stream.py :main.py
mpremote connect PORT reset
python3 host/check_rate.py --port PORT \
  --seconds 20 --save rate_deadline.csv
\end{lstlisting}
    \end{column}
  \end{columns}
  \small
  \begin{itemize}
    \item \textbf{You see:} \code{RESULT: PASS - the sampling rate is
          constant}.
    \item \textbf{You record:} the mean rate, the minimum period, and the
          maximum period.
    \item The check passes when the mean rate is within 1 percent of 50 Hz,
          at most 1 percent of the periods are outside 18 to 22 ms, and no
          line is lost.
  \end{itemize}
\end{frame}

\begin{frame}[fragile]{Part B: the same samples over Wi-Fi}
  \small
  \begin{enumerate}
    \item Connect the antenna to the XIAO.
    \item Write the name of the lab Wi-Fi, its password, and the IP address
          of your laptop in \code{board/config.py}.
    \item Copy the settings and start the Wi-Fi script. In a second
          terminal, measure the rate.
  \end{enumerate}
\begin{lstlisting}[style=shell]
mpremote connect PORT fs cp board/config.py :config.py
mpremote connect PORT run board/imu_wifi.py

python3 host/check_rate.py --udp 5005 --seconds 20     # second terminal
\end{lstlisting}
  \begin{itemize}
    \item The board sends the lines of 5 samples in one UDP packet: 10
          packets each second.
    \item UDP does not repeat a packet that the network loses. The sample
          number shows each lost sample.
    \item \textbf{You record:} the mean rate and the number of lost lines.
    \item \textbf{Question:} which connection loses samples, the serial
          port or Wi-Fi? Why?
  \end{itemize}
\end{frame}

% ------------------------------------------------------------------- Part C
\begin{frame}[fragile]{Part C: the logger (40 min)}
  \small
  Open \code{host/logger.py}. Complete the two functions with the mark
  \code{TODO (student)}.
  \vskip 0.3em
  \renewcommand{\arraystretch}{1.25}
  \begingroup\centering
  \begin{tabular}{@{}L{0.10\textwidth}L{0.22\textwidth}L{0.58\textwidth}@{}}
    \toprule
    \textbf{Task} & \textbf{Function} & \textbf{Example} \\
    \midrule
    C1 & \code{file\_name} & label \code{lift}, session \code{s2}, number
         7 gives \code{lift.s2.07.csv} \\
    C2 & \code{count\_lost} & the sample numbers 4, 5, 6, 9, 10 give 2 lost
         samples \\
    \bottomrule
  \end{tabular}\par\endgroup
  \vskip 0.3em
  Test the two functions with no board:
\begin{lstlisting}[style=shell]
python3 host/logger.py --self-test
\end{lstlisting}
\begin{lstlisting}
Task C1 (file_name):  complete
Task C2 (count_lost): complete
\end{lstlisting}
  \begin{itemize}
    \item The file name holds the label and the session. The notebook and
          Edge Impulse Studio read the label from the file name.
    \item The sample number of the board shows a lost sample. A recording
          with a lost sample is not a recording at 50 Hz.
  \end{itemize}
\end{frame}

\begin{frame}[fragile]{Part C: the protocol and the recordings}
  \small
  Write the protocol in the report before the first recording: one sentence
  for each class, and the person of each session.
  \vskip 0.3em
  \begin{columns}[T]
    \begin{column}{0.47\textwidth}
      \begin{itemize}
        \item \code{idle}: the kit lies on the table.
        \item \code{terrestrial}: left and right in a horizontal line.
        \item \code{lift}: up and down.
        \item \code{maritime}: slowly on all axes, as a boat on waves.
      \end{itemize}
    \end{column}
    \begin{column}{0.49\textwidth}
      \begin{itemize}
        \item Three sessions: \code{s1}, \code{s2}, \code{s3}.
        \item Change the person, the speed, or the size of the motion
              between the sessions.
        \item 4 recordings of 10 s for each class in each session.
      \end{itemize}
    \end{column}
  \end{columns}
\begin{lstlisting}[style=shell]
python3 host/logger.py --port PORT --label lift --session s1 --person A \
  --count 4 --plot
\end{lstlisting}
\begin{lstlisting}
Recording 1 of 4, class 'lift'. Start the motion now.
  recording 10 s ...
  saved data/lift.s1.01.csv: 500 samples, ... Hz, ... lost, ... clipped
  plot: data/plots/lift.s1.01.png
\end{lstlisting}
  Look at each plot. If a plot does not agree with its label, delete the
  file and record again.
  \source{The output shows the form of the lines.}
\end{frame}

\begin{frame}{Part C: the notebook has five tasks}
  \small
  \renewcommand{\arraystretch}{1.25}
  \begingroup\centering
  \begin{tabular}{@{}L{0.09\textwidth}L{0.36\textwidth}L{0.45\textwidth}@{}}
    \toprule
    \textbf{Task} & \textbf{You write} & \textbf{The notebook checks} \\
    \midrule
    1 & The mean sampling rate of a recording
      & Two known cases: 50 Hz and 41.67 Hz \\
    2 & The split by session
      & No session is in the two sets. Each class is in the two sets. \\
    3 & The number of windows of a recording
      & 500 samples give 41 windows of 100 samples with a stride of 10 \\
    4 & Your prediction for the leakage experiment
      & Not checked. You compare it with the result. \\
    5 & The data card in \code{report.md}
      & Not checked. The instructor reads it. \\
    \bottomrule
  \end{tabular}\par\endgroup
  \vskip 0.5em
  \begin{itemize}
    \item Start the notebook with \code{jupyter lab dataset.ipynb}.
    \item The notebook writes \code{data/split.json}. Day 3 reads this
          file.
    \item If \code{data/} has no recordings, the notebook uses a fallback
          dataset with simulated signals. Use it only if your group has no
          recordings.
  \end{itemize}
\end{frame}

\begin{frame}[fragile]{Part C: the leakage experiment}
  \small
  \renewcommand{\arraystretch}{1.2}
  \begingroup\centering
  \begin{tabular}{@{}L{0.10\textwidth}L{0.40\textwidth}L{0.40\textwidth}@{}}
    \toprule
    & \textbf{Procedure A (wrong)} & \textbf{Procedure B (correct)} \\
    \midrule
    Step 1 & Cut all recordings into windows & Split the recordings by
             session \\
    Step 2 & Normalize with the statistics of all windows
           & Cut the windows in each set \\
    Step 3 & Mix the windows: 80 percent training, 20 percent test
           & Normalize with the statistics of the training set \\
    \bottomrule
  \end{tabular}\par\endgroup
  \vskip 0.3em
  \begin{exerciseblock}
    \textbf{Prediction (task 4).} Which procedure gives the higher test
    accuracy? Which two numbers do you expect? Write them before you run
    the experiment.
  \end{exerciseblock}
\begin{lstlisting}
Procedure A, random split of the windows:    ... percent
Procedure B, test session s1                 ... percent
Procedure B, test session s2                 ... percent
Procedure B, test session s3                 ... percent
\end{lstlisting}
  The classifier is the same in the two procedures: the nearest neighbour
  on six features of a window. Only the split is different.
\end{frame}

% ------------------------------------------------------------------- Part D
\begin{frame}[fragile]{Part D: Edge Impulse (20 min)}
  \small
\begin{lstlisting}[style=shell]
python3 host/to_edge_impulse.py --data data
edge-impulse-uploader --category training ei_upload/training/*.csv
edge-impulse-uploader --category testing ei_upload/testing/*.csv
\end{lstlisting}
  \begin{enumerate}
    \item Convert: the time in milliseconds, the acceleration in m/s$^2$,
          one folder for each set.
    \item Make the project \code{motion-gNN} in Edge Impulse Studio, and
          upload. With no CLI, upload in the browser: \code{Data
          acquisition}, then \code{Add data}.
    \item Do not use the automatic split of the Studio. It does not know
          your sessions.
    \item Inspect: the four labels, one plot of each class, the duration
          of the training data and of the test data.
  \end{enumerate}
  \begin{block}{Before you leave: the Arduino firmware}
    In the Arduino IDE, set \code{Tools > Erase All Flash Before Sketch
    Upload > Enabled}. Upload \code{Labs/day01/sketches/blink/blink.ino}.
    Then set the option to \code{Disabled} again. If the upload fails,
    start bootloader mode first.
  \end{block}
\end{frame}
